\section{Results}
\label{sec:results}

We test four claims of the method in turn: that the calibration loss
has a sloppy spectrum (\S\ref{sec:results:bh-spectrum}); that the
diagnostic is usable from the first iterate
(\S\ref{sec:results:live}); that it survives non-MMD falsification
(\S\ref{sec:results:falsification}); and that it predicts where Adam
fails (\S\ref{sec:results:predicts-adam}).
Section~\ref{sec:results:network-sir} reports the same diagnostic
structure under Gumbel-Sigmoid surrogate gradients on a discrete-state
network model. Source code, seeds, and the figure-rendering scripts
for every experiment are released at \texttt{<repo-url-here>}
(anonymised at submission if double-blind).

\subsection{The Brock--Hommes calibration loss has a sloppy spectrum}
\label{sec:results:bh-spectrum}

The MMD-calibration loss of the canonical Brock--Hommes asset-pricing
model \citep{brock_hommes_1997,brock_hommes_1998} has an OPG spectrum
spanning \textbf{8.4 orders of magnitude} with condition number
$\lambda_1 / \lambda_P = 2.4 \times 10^8$ at a near-truth evaluation
point (Figure~\ref{fig:bh-spectrum}). The sloppiest direction $v_P$ is
the parameter $\beta$ alone --- the intensity of choice in the
switching dynamics --- which the MMD loss does not constrain. The
stiffest direction $v_1$ is the symmetric-bias combination
$(b_1 + b_2)/\sqrt{2}$, which no per-parameter sensitivity analysis
can surface but the eigendecomposition reads off directly. The data
does not constrain individual parameters; it constrains combinations.

\begin{figure}[t]
  \centering
  \includegraphics[width=\linewidth]{figures/fig1_spectrum.png}
  \caption{OPG geometry of the Brock--Hommes calibration loss. (a)~Per-seed gradient cloud in the top-2 OPG axes with the $1\sigma$ ellipse overlaid; the elongation \emph{is} the sloppiness. (b)~Eigenvalue spectrum with 95\% bootstrap confidence intervals; \textbf{8.4 OOM} span, condition $2.4 \times 10^8$. (c)~Parameter content $|V|$ of each eigendirection: $v_P$ is $\beta$ alone, $v_1$ is the symmetric-bias combination $(b_1 + b_2)/\sqrt{2}$.}
  \label{fig:bh-spectrum}
\end{figure}

\subsection{The diagnostic is live from the first iterate}
\label{sec:results:live}

The eigenstructure of $\widehat F$ is stable along the calibration
trajectory, not just at convergence. We bootstrap the top and bottom
eigenvectors at intermediate iterates of a far-from-equilibrium
calibration ($\|\theta_0 - \theta^*\| = 0.14$, $M=64$). The stiff
direction $v_1(\theta_t)$ stays within \textbf{$9.7^\circ$} of
$v_1(\theta^*)$ at $t=0$, narrowing to $1.9^\circ$ at convergence; the
sloppy direction $v_P(\theta_t)$ is within $20.3^\circ$ at $t=0$ and
within $1.1^\circ$ at convergence (Figure~\ref{fig:supp-trajectory} in
the appendix). The diagnostic is therefore usable from the first OPG
matrix the calibrator computes: the directions along which a
first-order optimiser will wander are visible before any compute has
been spent on the optimisation.

\subsection{Non-circular validation across three models}
\label{sec:results:falsification}

If the OPG eigenvectors identified an MMD-kernel artefact rather than
identifiability of the model, the discrepancies of
Equation~\eqref{eq:disc} along $v_1$ and $v_P$ would be comparable.
They are not. We measure stiff/sloppy ratios on three structurally
different agent-based models --- Brock--Hommes financial
\citep{brock_hommes_1997}, mean-field SIR
\citep{kermack_mckendrick_1927}, and a discrete-state network SIR with
Gumbel-Sigmoid surrogate gradients
\citep{maddison_concrete_2017,jang_gumbel_2017} --- against three
discrepancies the calibrator never saw: the first four moments, the
autocorrelation up to lag~20, and four tail quantiles. Every channel
on every model exceeds a $10^2$ ratio
(Figure~\ref{fig:falsification}). Brock--Hommes returns
$489\times / 581\times / 19{,}500\times$ for moments / ACF / tail
quantiles; mean-field SIR returns
$354{,}000\times / 930{,}000\times / 395{,}000\times$; network-SIR
returns $171{,}000\times / 238{,}000\times / 36{,}300\times$. The MMD
kernel mean embedding does not appear in the test; the result is
therefore not a self-fulfilling consequence of the calibration loss.

\begin{figure}[t]
  \centering
  \includegraphics[width=\linewidth]{figures/fig2_falsification.png}
  \caption{Non-circular validation of the OPG diagnostic. Three models (rows) $\times$ three discrepancies unrelated to the calibration kernel (columns: first four moments; autocorrelation up to lag~20; four tail quantiles) under the same $\alpha = 10^{-2}$ perturbation along the stiff $v_1$ and sloppy $v_P$ directions. Stiff/sloppy ratios are annotated per panel; aggregate ratios per model in the rightmost column. All twelve panels exceed $10^2$.}
  \label{fig:falsification}
\end{figure}

\subsection{The diagnostic predicts where Adam fails}
\label{sec:results:predicts-adam}

Decomposing the recovery error $\theta_T - \theta^*$ along the OPG
eigenbasis at $\theta^*$, the median squared component tracks the
inverse eigenvalue scaling: small-$\lambda_k$ directions are the
directions along which every optimiser fails
(Figure~\ref{fig:predicts-adam}). Adam fails along the sloppy
direction more than OPG-preconditioned Levenberg-Marquardt does: at
medium difficulty ($\|\theta_0 - \theta^*\| = 0.15$, $5$ random
unit-vector inits), the cross-optimiser ratio
Adam~$v_P$~/~OPG~$v_P$ is \textbf{$34\times$} and
Adam~$v_P$~/~OPG~$v_1$ is \textbf{$9.4 \times 10^4$}. Across 15
Brock--Hommes seeds and 10 SIR seeds, Adam diverges from $\theta^*$ in
$25/25$ runs (95\% LCB $\approx 88\%$). The diagnostic flags the
failure direction before the failure itself need be observed.

\begin{figure}[t]
  \centering
  \includegraphics[width=\linewidth]{figures/fig3_predicts_adam.png}
  \caption{Recovery error decomposed along the OPG eigenbasis at $\theta^*$. (a)~Median (IQR) squared component per eigendirection over $5$ random unit-vector inits at distance $0.15$ from $\theta^*$; Adam's $v_5$ bar towers above OPG's and SGD's. (b)~Median squared error against eigenvalue; small-$\lambda$ directions are the failure directions for every optimiser, with Adam's failure systematically larger. Cross-optimiser ratios: Adam~$v_P$/OPG~$v_P = 34\times$; Adam~$v_P$/OPG~$v_1 = 9.4 \times 10^4$.}
  \label{fig:predicts-adam}
\end{figure}

\subsection{The diagnostic survives discrete states under Gumbel-Sigmoid surrogates}
\label{sec:results:network-sir}

The standing concern with differentiable agent-based modelling is that
gradients computed through a discrete-state simulator with
Gumbel-Sigmoid surrogate transitions
\citep{maddison_concrete_2017,jang_gumbel_2017} are biased by
construction, and could mislead any identifiability diagnostic that
consumes them. We test this directly on a network-SIR model with
per-node Gumbel-Sigmoid infection and recovery transitions
($N=250$, mean degree $6$, $\tau_{\mathrm{Gumbel}}=0.5$). The OPG
eigenvectors recover the same qualitative stiff/sloppy structure as
the smooth mean-field SIR: the stiff direction is dominated by the
initial-infection fraction with a contact-rate correction,
$v_1 \approx 0.89\, I_0 + 0.46\, \beta$; the sloppy direction is the
lockdown intensity, $v_P \approx 0.995\, f_{\mathrm{lock}}$ --- the
early-weak-versus-late-strong policy degeneracy a public-health
analyst would derive from the model equations
(Figure~\ref{fig:network-sir}). The spectrum widens to
\textbf{20.7 orders of magnitude} and the non-MMD falsification ratios
reach $10^4$--$10^5\times$ even under the surrogate bias, so the
biased gradients identify real identifiability structure rather than
surrogate artefacts.

\begin{figure}[t]
  \centering
  \includegraphics[width=\linewidth]{figures/fig4_network_sir.png}
  \caption{Network-SIR diagnostic with Gumbel-Sigmoid surrogate gradients. (a)~Daily new infections at $\theta^*$ (with lockdown) versus a no-lockdown counterfactual. (b)~OPG spectrum with 95\% bootstrap confidence intervals; \textbf{20.7 OOM} span. (c)~Eigenvector content; $v_1 \approx 0.89\, I_0 + 0.46\, \beta$, $v_P \approx 0.995\, f_{\mathrm{lock}}$. (d)~Mean-field versus network-SIR spectra: the qualitative structure transfers under surrogate-gradient bias.}
  \label{fig:network-sir}
\end{figure}
